\chapter{Background}
\section{Introduction to Explainability for LLMs}
\subsection{Definition and Purpose}
Large language models (LLMs) have seen increasing adoption in recent years for task such as question answering \cite{allemangIncreasingAccuracyLLM2025}, writing \cite{liangWidespreadAdoptionLarge2025} and software development \cite{zhangSurveyLargeLanguage2026} but also in high-stakes domains such as healthcare \cite{maityLargeLanguageModels2025} or law \cite{laiLargeLanguageModels2024} \cite{hungWalkingTightropeEvaluating2023}. However, since LLMs are complex deep learning models with millions of neurons and billions of parameters, they are considered black boxes  \cite{qamarUnderstandingBlackboxInterpretable2023}. This makes it challenging to comprehend and evaluate decisions.

Explainable AI (XAI) is a research field that seeks to make artificial intelligence algorithms transparent and understandable \cite{hsiehComprehensiveGuideExplainable2024} - making the inherent black box less opaque. Central to this research field is the concept of Explainability whose clear definition is still disputed \cite{roscherExplainableMachineLearning2020}. This term is described as the "level of understanding how the AI-based system came up with a given result" in the pertaining ISO-Standard \cite{ISOIECTR}. Explainability attempts to justify decisions, showing that they are correct and non-biased to build trust with users. \cite{aliExplainableArtificialIntelligence2023}.

The term "Interpretability" is often used in a similar context even though there is a fine distinction. Interpretability is about understanding the inner-workings of an AI model where Explainability focuses on making the decision-process retraceable \cite{garouaniInvestigatingDualityInterpretability2024} \cite{aliExplainableArtificialIntelligence2023}.

\subsection{Transparency as Regulatory Requirement}
\label{sec:transparency_as_regulatory_requirement}
The recent surge in importance of the XAI-research field has also been fueled by new regulatory measures. While multiple countries have introduced their own legislative frameworks to govern the use of AI such as China \cite{zouNavigatingChinasRegulatory2025} and the US \cite{BlueprintAIBill}, the European Union's AI Act \cite{RegulationEU20242024} can be seen as the most pertinent to explainability. Within this law, individuals impacted by high-risk AI decisions have the right to obtain clear and meaningful explanations of the AI’s role and the main contributing factors to the decision made. Furthermore high-risk AI systems must offer transparency and be supervisable under this legislation \cite{RegulationEU20242024}. The "Right of Explaination" from the General Data Protection Regulation (GDPR) is also of relevance to AI systems within the EU. \cite{nisevicExplainableAICan2024}

\subsection{Audience Groups}
Providing explanations for reasoning or decisions made by an AI system can increase trust in the system with end-users and experts \cite{viloneNotionsExplainabilityEvaluation2021} \cite{herreraMakingSenseUnsensible2025} and can assist researchers and developers in debugging and fine-tuning the model \cite{szymanskiDisentanglingStakeholderRole2025}. Moreover as argued in chapter \ref{sec:transparency_as_regulatory_requirement}, explainability can be required by regulations and auditors forming another stakeholder group for XAI systems \cite{herreraMakingSenseUnsensible2025}  \cite{szymanskiDisentanglingStakeholderRole2025}.

It has also been noted that different audiences perceive and consume explanations differently \cite{herreraReflectionsAttentivenessEXplainable2025}\cite{herreraMakingSenseUnsensible2025}. For instance, where developers would expect technical data about the inner-workings of the system, end-users might only benefit from easy-to-understand explanations in natural language. Adapting explanations given by an Artificial Intelligence based on the expectations of the recipient has been subject to recent research \cite{hoffmanExplainableAIRoles2023} \cite{mavrepisXAIAllCan2024}.

\subsection{Aspects and Dimensions}
\label{sec:aspects_and_dimensions}
Herrera \cite{herreraMakingSenseUnsensible2025} points out that whereas traditional AI models usually follow fixed deterministic rules for generating an explanation, LLMs complicate this process with their more complex, advanced features. The author therefore suggests four dimensions of explainability for large language models:
\begin{itemize}
\item \textbf{Faithfulness}: Faithfulness describes if the explanation does correspond to the internal reasoning of the model \cite{jacoviFaithfullyInterpretableNLP2020}.
\item \textbf{Truthfulness}: Truthfulness describes if the explanation does conform to real-world facts drawn from an external authoritative reference. A truthful explanation does not contain factual errors, omissions or distortions \cite{herreraMakingSenseUnsensible2025}.
\item \textbf{Plausibility}: Plausibility describes the degree, to which an explanation is convincing and sounds reasonable to humans. \cite{zhuExplanationEraLarge2024}\cite{jacoviFaithfullyInterpretableNLP2020}
\item \textbf{Contrastivity}: Contrasitivy measures how well the explanation highlights differences between competitive pieces of information. In classification tasks, an XAI system with high contrastivity clearly outlines \textit{why} it chose one class over another. \cite{herreraMakingSenseUnsensible2025} \cite{mershaUnifiedFrameworkNovel2025} \cite{stepinSurveyContrastiveCounterfactual2021}.
\end{itemize}

It must be noted that there is not wide scientific consence within the XAI community yet on which aspects are essential for an explanation to be considered good since descriptions, names and definitions of these factors can vary considerably. \cite{phillipsFourPrinciplesExplainable2021} \cite{aliExplainableArtificialIntelligence2023} \cite{lopesXAISystemsEvaluation2022} \cite{mershaUnifiedFrameworkNovel2025} This effect is amplified by regulatory and standardization frameworks that remain rather vague in this respect \cite{RegulationEU20242024} \cite{ISOIECTR}. Nevertheless, Herreras \cite{herreraMakingSenseUnsensible2025} desiderata of Faithfulness and Plausibility in particular are supported by other literature such as \cite{jacoviFaithfullyInterpretableNLP2020}, \cite{zhuExplanationEraLarge2024} and \cite{turpinLanguageModelsDont2023}. Others like Lopes et al. \cite{lopesXAISystemsEvaluation2022} refer to Faithfulness as "Fidelity" and split Plausibility into other human-centered sub-aspects.

\subsection{Challenges and Opportunities in LLM explainability research}
\label{sec:challenges_in_llm_explainability_research}
According to Zhao et al. \cite{zhaoExplainabilityLargeLanguage2024} explainity research in LLMs is difficult because understanding their processes becomes inviable due to their vast size and complex structure. Traditional methods like gradient-based attribution \cite{sundararajanAxiomaticAttributionDeep2017} and SHAP values \cite{lundbergUnifiedApproachInterpreting2017} are impractical for state-of-the-art LLMs with billions of parameters. This complexity also hinders in-depth analysis, debugging, and diagnosing models. Moreover, the intransparent black-box nature of LLMs makes amplifies this problem \cite{qamarUnderstandingBlackboxInterpretable2023}. Many mainstream commercial models are only available via an API \cite{mattonWalkTalkMeasuring2024}, where the weights are not released to users or researchers.

Despite of this, several technical approaches have been developed to take a glimpse into the black-box and retrace important pathways within an LLM for a decision. These can for brevity not be discussed in this thesis, but the most prominent are presented in \cite{zhaoExplainabilityLargeLanguage2024} and \cite{luoUnderstandingUtilizationSurvey2024}.

Still, since LLMs are designed to handle Natural-Language-Processing Tasks (NLP) well, there is also a great opportunity in letting the model give explanations in natural human language  \cite{zhaoExplainabilityLargeLanguage2024}. This, most notably, includes Chain-of-Thought (CoT) explanations, where LLMs reconstruct their own thought process. Instructing the LLM to generate such CoT-explanations has furthermore been shown to increase the performance in complex reasoning tasks \cite{weiChainofthoughtPromptingElicits2022}. From now on in this work, CoT shall be the method used to bring about an explantaion from the Large Language model.

\section{Faithfulness of natural language-explanations for LLMs}
\subsection{Definition}
Faithfulness describes the degree to which the explanation given by an AI-powered system is grounded in the actual internal reasoning \cite{herreraMakingSenseUnsensible2025} \cite{jacoviFaithfullyInterpretableNLP2020} \cite{ribeiroWhyShouldTrust2016}. It requires explanations to reflect the model's true decision path, not merely offer rationalizations. \cite{herreraMakingSenseUnsensible2025}

\subsection{Significance and Motivation for Evaluation}
The desiderata of explanations (\ref{sec:aspects_and_dimensions})  are independent of each other: It's possible for an XAI system to satisfy one criteria but not another. \cite{herreraMakingSenseUnsensible2025} \cite{jacoviFaithfullyInterpretableNLP2020} An LLM might output an natural language explanation sounding convincing to users (high Plausability) but purely made up and not grounded in reality (low Faithfulness). This scenario can pose significant risks and lead to misplaced trust in those systems. \cite{agarwalFaithfulnessVsPlausibility2024} \cite{madsenAreSelfexplanationsLarge2024} \cite{herreraMakingSenseUnsensible2025}

In particular, research by Turpin et al. \cite{turpinLanguageModelsDont2023} showed that LLMs can systematically hide internalized biases and offer plausible yet unfaithful justifications. When influenced to make use of prejudices in their experiment, models usually used social biases such as gender or race as primary decision factors, but rarely mentioned them, instead resorting to conduct or circumstances for their justification. This can be particular concerning in use cases such as recidivism prediction, where fairness is undermined by this behavior. \cite{farayolaFairnessAIPredicting2023} \cite{jacoviFaithfullyInterpretableNLP2020} Therefore, not only verifying that the explanation is reasonable but also showing that the given explanation does adhere to the actual thought process is of uttermost importance, particularly for LLMs in critical domains.

\subsection{Estimation and Quantification Approaches}
As touched upon in chapter \ref{sec:challenges_in_llm_explainability_research} the complex, intransparent nature of Large Language Models poses unique challenges. This challenge is particular pronounced in assessing the faithfulness of CoT-explanations. Where for instance Plausibility can be tested with humans \cite{herreraMakingSenseUnsensible2025}, faithfulness is inherently tied to the internals of a model \cite{parcalabescuMeasuringFaithfulnessSelfconsistency2024}.

One of the first works in regards to estimating faithfulness of LLMs is by Turpin et al \cite{turpinLanguageModelsDont2023}. The paper contributes - besides showing that CoT-explanations can systematically skip contributing factors for their decision - a quantification of unfaithfulness on the BBH-Dataset by measuring the drop in model accuracy when biasing features are added to inputs that the model fails to mention in its explanations. Their approach is labor-intensive: CoT-explanations generated by the LLM were checked by humans if the biased feature are present. According to the authors, their quantification method only evaluates faithfulness over smaller changes to the input where a practically usable metric could dig deeper.

Another paper by Chen et. al. \cite{chenModelsExplainThemselves2024} introduces the concept of counterfactual simulatability. This is described as the degree to which humans can infer "the model’s outputs on diverse counterfactuals of the explained input", meaning a human is able to predict the LLM to answer to closely-related yes or no questions in a expected manner after reading the explanation to an initial question. Based on this concept, the researchers propose two faithfulness metrics: simulation generality calculates the "diversity of the counterfactuals relevant to the explanation" and simulation precision measures the fraction of counterfactual questions to which human expectations line up with the model's output. Their evaluation pipeline first pulls a question from the dataset and gives it the LLM to generate an answer and create an explanation. This explanation is then used to generate pertaining counterfactuals. This task however is done by an LLM for better scaling with large datasets. A human then builds a mental model based on the original explanation and states his expectations for the answers to the counterfactual questions. The counterfactuals are then used as input to the LLM. Finally, generality and precision scores are calculated for the given explanation. This approach notably works both with ad-hoc and Chain-of-Thought explanations and is an early example of using auxiliary LLMs to generate counterfactuals, inheriting the limitations that come with a LLM-as-a-Judge choice \cite{mattonWalkTalkMeasuring2024}.

Another example of using another AI to generate counterfactuals is from Atanasova et al. \cite{atanasovaFaithfulnessTestsNatural2023}. The researchers can be credited with setting the cornerstone for a line of work sometimes called "counterfactual edits" \cite{parcalabescuMeasuringFaithfulnessSelfconsistency2024} \cite{mattonWalkTalkMeasuring2024} also used by later metrics \cite{siegelProbabilitiesAlsoMatter2024} \cite{mattonWalkTalkMeasuring2024} where the LLMs response to the counterfactual is used as basis for the faithfulness estimation . For their paper, the authors trained an auxillary model to add words to the input, giving yield to counterfactuals. The original input and these counterfactuals are then being given to the LLM which should be evaluated. The two responses and the explanation to the counterfactual input are then evaluated: If the prediction changes with the counterfactual, and the inserted words are not mentioned in the explanation, the explanation is seen as unfaithful. The benifit is, that this LLM-based evaluation method is without the need of human annotators.  \cite{parcalabescuMeasuringFaithfulnessSelfconsistency2024} There are still limitations of this approach, for instance the input modifications might cause the model to focus elsewhere, meaning its prediction isn't necessarily derived from the specific insertion. \cite{atanasovaFaithfulnessTestsNatural2023}. But Siegel et. al. \cite{siegelProbabilitiesAlsoMatter2024} also see two more practical considerations with Atanasova et al's implementation of "counterfactual edits": 1. The test does not investigate whether high-impact features are more likely to be mentioned than low-impact ones. An AI system might, as argued by Siegel et. al\cite{siegelProbabilitiesAlsoMatter2024}, reach 100% faithfulness by just repeating the entire input text. 2. The Binary impact measurement only cares if a prediction probability of a class reaches 50% for an intervention, missing larger probability shifts (e.g. from 6% to 45%).

In order to account for these drawbacks, Siegel et. al. \cite{siegelProbabilitiesAlsoMatter2024} significantly improve Atanasova et al.'s \cite{atanasovaFaithfulnessTestsNatural2023} faithfulness metric with their Correlational Counterfactual Test (CCT). In core they measure faithfulness as the "correlation between correlation between intervention impact scores and explanation mention scores" \cite{mattonWalkTalkMeasuring2024}. First, the LLM is being fed an unmodified question from a dataset. Then the input is perturbed and given to the LLM, obtaining the LLMs responses as probability distribution (impact values). They inherit the counterfactual creation procedure mostly from the original work, while also using an auxillary LLM to filter out non-sensical interventions. Then, the Total Variation Distance (TVD) is being calculated, measuring the shift between original predictions and predictions on counterfactuals. Then, the explanation given by the LLM is being checked by substring matching if the explanation mentions the intervention, giving a binary value. Then this procedure is repeated for multiple questions from a data set. Finally, to return the CCT the Pearson correlation coefficient between the impact values of the intervention and the explanation mentions is computed. By using Atanasova et al. \cite{atanasovaFaithfulnessTestsNatural2023} approach as basis, some limitations are inherited, particularly in regards to the counterfactual creation. For instance, counterfactuals only have a single adjective or adverb added, leading to less creative counterfactuals. Furthermore, the score does not account for the explanation to mention a synonym but not the added input word itself.

Drawing from the improvements made on counterfactual edits, Matton et. al. \cite{mattonWalkTalkMeasuring2024} propose their "Causal Concept Faithfulness Metric". Unlike previous, they consider concepts such as "gender", "race" or "time of day" instead of low-level tokes or words. Furthermore, their method does not only provide a quantitive score but also exposes semantic patterns in regards to faithfulness. They treat faithfulness as the alignment between truly influential concepts in a models prediction and the importance the concepts have according to the models (in their explanation). They use an auxillary LLM to either replace the value of a concept or to remove the concept from the input. Then like in CCT-Metric, the original and counterfactual questions are being given to the LLM and prediction distributions are collected, resulting in a true causal effect. Now they extract which concepts are being deemed important in the explanation using another helper LLM which then can be used for calculating the Explanation-implied effect. The Pearson correlation between the true causal effects and the Explanation-implied effects then results in a faithfulness score. One limitation comes with the LLM-automated counterfactual creation and explanation extraction which might result in errors and poses a risk of circularity in the evaluation process. Furthermore, this metric does not work with concepts that are correlated. For brevity, only a short overview ofer the Causal Concept Faithfulness Metric is given here, however since it is central to this work, it will be discussed in more technical detail in chapter \ref{sec:causal_concept_metric}.

Paracalabescu and Frank \cite{parcalabescuMeasuringFaithfulnessSelfconsistency2024} argue that current metrics measure self-consistency of outputs rather than true faithfulness. They then present their own self-consistency metric called CC-SHAP which does not depend on modifying the input. As rationale they compare the degree to which a input contributes to the answer and the explanation on a token level. By measuring the convergence between these using Shapley-Values a continuous result is obtained. Their approach shows strength in not only not needing input perturbations but also does not need semantic evaluation of explanations, works for both CoT and ad-hoc explanations, show patterns of (in-)consistency on a token level and do not need a AI-model anywhere in the evaluation process. However, these benefits come with a computational overhead that can be more pronounced than other methods.

Lanham et al. \cite{lanhamMeasuringFaithfulnessChainofThought2023} has experimented with corrupting the CoT in order to investigate certain hypothesis about CoT-reasoning. First, the model is asked to reason step for step about a multiple choice problem. Given the original CoT, the LLM is prompted to give a final answer. The CoT is then edited. For this they either trunkate the CoT, adding mistakes to it (via auxillary model), paraphrase it (by rewording the beginning of the CoT and let the model regenerate the rest) as well as use filler tokens. Then the prediction is collected again for the edited CoT. If the two predictions match despite the intervention, this is seen as unfaithfulness. Since their metric is primarily only used as the means to collect evidence about the research hypothesis, it is not entirely clear how it compares to other metrics. Paracalabescu and Frank \cite{parcalabescuMeasuringFaithfulnessSelfconsistency2024} however argue, that this approach has the requirement that the LLM absolutely needs the CoT to give the correct prediction. However, Lanham et al. \cite{lanhamMeasuringFaithfulnessChainofThought2023} own tests have shown that models with CoT only marginally improve their accuracy, compared to those who don't have a CoT-explanation \cite{parcalabescuMeasuringFaithfulnessSelfconsistency2024}, putting the practical usefulness of this method into question.

One last approach shall be presented: A recently published work by Tutek et al. \cite{tutekMeasuringChainThought2025} tests an entirely novel approach called Faithfulness by Unlearning Reasoning steps (FUR). They propose if a reasoning step is truly used by the model to arrive at its answer, then removing the information from the model’s parameters should change its prediction. Namely, for a given LLM $M$ and a test question, they let $M$ generate a CoT-reasoning. In the next step, the explanation received is split into specific steps. Then for each reasoning step they unlearn the knowledge from $M$'s parameters. This gives yield to a fine-tuned model $M'$ which is inhibited to predicting the words of the reasoning step in this context. Then they ask the original model $M$ and the modified version $M'$ the same question, now without asking for a CoT-explanation and compare their chosen results. For their ff-hard metric, measuring the faithfulness of the entire CoT, a binary result is returned- A result of 1 indicates that both models disagree in their most likely answer - showing high faithfulness, 0 if the chosen answer still remains the same and the given explanation does likely not reflect the actual reasoning. One major strengh is that this method is void of the downfalls of context-pertupation methods like the model being able to reconstruct the context. The authors themselves however outline that high recall cannot be guaranteed: if unlearning a step does not change the prediction, it could be because unlearning failed, because multiple reasoning paths exist, or because the step was genuinely not used. The FUR-approach of course is also not applicable for closed-source API-only models.



\section{The Causal Concept Faithfulness Metric}
\label{sec:causal_concept_metric}

\section{The Correlated Concepts Problem}

